\documentclass[12pt]{article}

\usepackage{amssymb} % symbol kąta
\usepackage[polish]{babel} % polskie nazwy
\usepackage[T1]{fontenc} % polskie znaki
\usepackage[margin=1.0in]{geometry} % marginesy
\usepackage[utf8]{inputenc}
\usepackage{listingsutf8} % bloki kodu
\usepackage{color} % kolory
\usepackage{indentfirst} % wcięcie w pierwszej linii paragrafu
\usepackage{graphicx} % obrazy
\usepackage{float} % dla image [H]
\usepackage{fancyhdr} % nagłówki stron

% bez odstępu w itemize
\let\tempone\itemize
\let\temptwo\enditemize
\renewenvironment{itemize}{\tempone\setlength{\itemsep}{0cm}}{\temptwo}

\definecolor{lbcolor}{rgb}{0.9,0.9,0.9}  
\graphicspath{ {img} }
% obrazek {nazwa.jpg}{opis}
\newcommand{\imgcustomsize}[3]{
	\begin{figure}[H]
		\centering
		\includegraphics[width=#3\textwidth]{#1}
		\caption{#2}
		\label{#1}
	\end{figure}
}
\newcommand{\img}[2]{\imgcustomsize{#1}{#2}{0.5}}

% dwa obrazki
\newcommand{\imgsidebyside}[4]{
	\begin{figure}[H]
		\centering
		\begin{minipage}{.5\textwidth}
			\centering
			\includegraphics[width=.8\linewidth]{#1}
			\caption{#2}
		\end{minipage}%
		\begin{minipage}{.5\textwidth}
			\centering
			\includegraphics[width=.8\linewidth]{#3}
			\caption{#4}
		\end{minipage}
	\end{figure}
}


% tabela z wynikami {nazwa.txt}{opis}
\newcommand{\tab}[2]{
\begin{table}[H]
	\centering
	\caption{#2}
	\vspace{0.3cm}
	\begin{tabular}{c|c|c|c}
		\lefthead{$x_1$}&\thead{$y_1$}&\thead{$x_2$}&\thead{$y_2$}\\
		\hline
		\input{#1}
	\end{tabular}
\end{table}
}

% nagłowek tabeli
\newcommand*{\thead}[1]{\multicolumn{1}{|c}{\bfseries #1}}
\newcommand*{\lefthead}[1]{\multicolumn{1}{c}{\bfseries #1}}

% nagłówki stron
\pagestyle{fancy}
\fancyhf{}
\rhead{Obliczanie grafu widoczności}
\lhead{Algorytmy geometryczne}
\rfoot{\thepage}

\begin{document}
	\lstset{backgroundcolor=\color{lbcolor}, language=Python, inputencoding=utf8}
	
	\title{\includegraphics{agh.jpg} \\ Obliczanie grafu widoczności \\
		\vspace{0.4cm}
		\large Algorytmy geometryczne \\
		\large Projekt}
	\author{Norbert Morawski\\Dariusz Piwowarski}
	
	\maketitle
	\pagebreak
	\tableofcontents
	\pagebreak
	
	\section{Informacje techniczne}
	% \hspace{0.5cm}
	\subsection{Platforma}
		\begin{tabular}{l|l}
			System&Arch Linux/Windows 10\\
			Architektura&\texttt{x86\_64}\\
			Język programowania&Python\\
			Wersja interpretera&3.9.7\\
		\end{tabular}

	\subsection{Użyte biblioteki i narzędzia}
		\begin{tabular}{l|l}
			NumPy&1.21.2\\
			MatPlotLib&3.4.3\\
			Jupyer Notebook&6.4.4
		\end{tabular}
	
	\subsection{Opis programu}
		% TODO może coś więcej tutaj
		Program został napisany w języku Python. Interfejs dla użytkownika dostępny jest w pliku \lstinline|.ipynb|, wykorzystuje on narzędzie Jupyter Notebook.
	
	\subsection{Moduły}
		Poniżej znajduje się opis modułów i metod w nich występujących.
	
		\paragraph{\lstinline|create_visibility_graph.py|}
		Zawiera implementację algorytmu obliczania grafu widzialności.
		\begin{itemize}
			\item \lstinline|find_internal_lines| -- Implementuje algorytm, który znajduje dla każdego punktu figury, inne wierzchołki tego samego wielokąta, takie że linia je łącząca jest skierowana do wnętrza figury. Złożoność $ O(n^2) $,
			\item \lstinline|create_visibility_graph| -- Tworzy graf  widzialności.
		\end{itemize}
	
	\paragraph{\lstinline|VisibilityGraph.py|}
		Klasa zawierająca reprezentację finalnego grafu widzialności.
		\begin{itemize}
			\item \lstinline|__init__| -- Konstruktor,
			\item \lstinline|add_edge| -- Dodaje krawędź do grafu,
			\item \lstinline|get_points| -- Zwraca listę krotek reprezentujących zbiór wierzchołków grafu,
			\item \lstinline|get_lines| -- Zwraca listę krawędzi,
			\item \lstinline|get_lines_separately| -- Zwraca listę list krawędzi z podziałem na wierzchołki początkowe krawędzi,
			\item \lstinline|__repr__| -- Zwraca reprezentację tekstową klasy.
		\end{itemize}
	
	\paragraph{\lstinline|Sweeper.py|}
		Klasa zawierająca aktualny stan miotły (punkt początkowy i końcowy)
	
	\paragraph{\lstinline|shortest_path.py|}
		Implementacja algorytmu Dijkstry, który wyznacza najkrótszą ścieżkę pomiędzy przeszkodami.
		\begin{itemize}
			\item \lstinline|dijkstra| -- Implementuje algorytm Dijkstry. Złożoność $ O(n \cdot log(n)) $,
			\item \lstinline|shortest_path| -- Funkcja wyznaczająca najkrótszą ścieżkę pomiędzy dwoma zadanymi punktami, na płaszczyźnie zawierającej przeszkody w postaci wielokątów.
		\end{itemize}
	
	\paragraph{\lstinline|Point.py|}
		Klasa zawierająca reprezentację punktu na płaszczyźnie.
		\begin{itemize}
			\item \lstinline|__init__| -- Konstruktor,
			\item \lstinline|with_line| -- Dodaje krawędź incydentną z tym wierzchołkiem,
			\item \lstinline|has_line_to| -- Sprawdza czy wierzchołek jest incydentny z podaną krawędzią,
			\item \lstinline|has_internal_line_to| -- Sprawdza czy krawędź pomiędzy tym a podanym wierzchołkiem jest skierowana do wnętrza figury, 
			\item \lstinline|__repr__| -- Zwraca reprezentację tekstową klasy.
		\end{itemize}
	
	\paragraph{\lstinline|Line.py|}
		Klasa zawierająca reprezentację odcinka na płaszczyźnie. Wewnętrznie przechowuje jego początek i koniec.
		\begin{itemize}
			\item \lstinline|__init__| -- Konstruktor,
			\item \lstinline|__repr__| -- Zwraca reprezentację tekstową klasy,
			\item \lstinline|__eq__| -- Sprawdza czy linie są sobie równe (porównuje czy mają takie same końce i początki),
			\item \lstinline|__gt__| -- Sprawdza czy linia następuje po podanej linii w strukturze stanu.
		\end{itemize}
	
	\paragraph{\lstinline|Figure.py|}
		Klasa zawierająca reprezentację figury na płaszczyźnie. Wewnętrznie przechowuje listę punktów należących do figury, w kolejności przeciwnej do ruchu wskazówek zegara.
		\begin{itemize}
			\item \lstinline|__init__| -- Konstruktor,
			\item \lstinline|__repr__| -- Zwraca reprezentację tekstową klasy.
		\end{itemize}
	
	\paragraph{\lstinline|plotter/Plotter.py|}
		Klasa obsługująca generowanie wykresów krok-po-kroku.
		\begin{itemize}
			\item \lstinline|__init__| -- Konstruktor,
			\item \lstinline|init_limits| -- Pobiera maksymalne i minimalne współrzędne punktów w grafie,
			\item \lstinline|new_partial_plot| -- Tworzy nowy częściowy wykres,
			\item \lstinline|partial_plot| -- Rysuje nowy częściowy wykres,
			\item \lstinline|sum_up| -- Rysuje podsumowanie jednej iteracji zewnętrznej pętli algorytmu.
		\end{itemize}
	
	\paragraph{\lstinline|plotter/SequencePlotter.py|}
		Klasa pomocnicza obsługująca generowanie osobnych wykresów w osobnych plikach, za każdym razem inkrementując numer wykresu.
		\begin{itemize}
			\item \lstinline|__init__| -- Konstruktor,
			\item \lstinline|next| -- Tworzy nowy wykres.
		\end{itemize}
	
	\paragraph{\lstinline|plotter/Plot.py|}
		Klasa pomocnicza zawierająca podstawowe funkcje obsługi wykresów.
	
	\paragraph{\lstinline|plot_tool.py|}
		Dostarczone narzędzie graficzne. Została dopisana funkcjonalność opisywania wierzchołków kolejnymi literami.
	
	\paragraph{\lstinline|plot_tool_helpers.py|}
		\begin{itemize}
			% TODO nie wiem
			\item \lstinline|visibility_graph_scenes| -- 
			\item \lstinline|visibility_graph_scenes_separately| -- 
			\item \lstinline|shortest_path_scene| -- 
			\item \lstinline|get_figures_from_plot| -- 
			\item \lstinline|get_points_from_plot| -- 
		\end{itemize}
	
	\paragraph{\lstinline|helpers.py|}
		Zawiera funkcje pomocnicze do obliczania grafu widoczności.
		\begin{itemize}
			\item \lstinline|dist_p_to_intersection(p, w, line)| -- Oblicza dla zadanych punktów $ p $, $ w $ i pewnego odcinka, odległość punktu $ p $ od przecięcia $ pw $ z tym odcinkiem. Jeśli są równoległe zwracany jest środek zadanego odcinka,\\
			\item \lstinline|intersection(line1, line2)| -- Zwraca $ True $/$ False $ w zależności, czy zadane odcinki się przecinają,\\
			\item \lstinline|orient(a, b, c)| -- Określa w jaki sposób zorientowane są trzy zadane punkty. Zwraca wartość $ 1 $, gdy są zorientowane przeciwnie do ruchu wskazówek zegara; $ 0 $, gdy są współliniowe; $ -1 $ gdy ich orientacja jest zgodna z ruchem wskazówek zegara,\\
			\item \lstinline|dist(a, b)| -- Oblicza odległość pomiędzy dwoma punktami na płaszczyźnie.
		\end{itemize}

	\section{Informacje użytkownika}
	
	\subsection{Uruchomienie programu}
		% TODO może coś więcej tutaj
		Należy w katalogu z plikiem \lstinline|project.ipynb| uruchomić narzędzie Jupyter Notebook.
	
	\subsection{Funkcje programu}
			% TODO Tutaj opis Jupytera
	
	\section{Sprawozdanie}
	
	\subsection{Opis problemu}
		\paragraph{Graf widzialności}
			Jest to graf gdzie zbiorem wierzchołków jest zbiór punktów na płaszczyźnie, a krawędzie istnieją tylko kiedy wierzchołki "widzą się", tzn. nie istnieje pomiędzy nimi przeszkoda w postaci figury.
		
		\paragraph{Wstęp}
			Trywialny algorytm wymagałby sprawdzenia każdej pary wierzchołków w czasie $ O(n^2) $. Każde takie sprawdzenie wymagałoby przejrzenia wszystkich innych wierzchołków w czasie $ O(n) $, co prowadziłoby do algorytmu o złożoności $ O(n^3) $. Jednak jeżeli wykorzystamy zrównoważone drzewo binarne do wyszukiwania przeszkód, jesteśmy w stanie obniżyć złożoność do poziomu $ O(n^2 \cdot log(n)) $.
		
		\paragraph{Algorytm}
			Problem sprowadza się do wyznaczenia widzianych wierzchołków dla każdego z $ n $ punktów. Do tego celu wykorzystamy algorytm zamiatania. Płaszczyzna jest zamiatana poprzez obrót miotły zaczepionej w rozważanym punkcie, zgodny z ruchem wskazówek zegara.
			
		\paragraph{Zdarzenia i stan miotły}
		Zdarzeniami będą punkty na płaszczyźnie, a struktura stanu będzie przechowywać aktualnie przecinane odcinki w kolejność od najbliższego punktowi zaczepienia do najdalszego.
	
	\subsection{Algorytm wyznaczania widocznych wierzchołków dla punktu $ p $} \label{sssec:algo}
		Algorytm polega na zamiataniu płaszczyzny i odbywa się wg następujących kroków:
		\begin{enumerate}
			\item Posortuj wierzchołki na podstawie kąta (zgodnego z ruchem wskazówek zegara) jaki półprosta z p do danego wierzchołka tworzy z dodatnią półosią OX. Gdy kąt jest taki sam, wierzchołek bliższy p powinien znajdować się przed wierzchołkiem dalszym. Wierzchołki w tej kolejności umieść w strukturze zdarzeń.
			\item Zaczep miotłę w badanym punkcie i zwróć ją w kierunku dodatniej półosi OX.
			\item Dodaj do struktury stanu krawędzie przecinające miotłę w położeniu początkowym.
			\item Dla kolejnych punktów $ w $ z struktury zdarzeń:
			\begin{itemize}
				\item Ustaw miotłę jako półprostą z $ p $ do $ w $.
				\item Sprawdź czy punkt $ w $ jest widoczny, jeśli tak dodaj odpowiednią krawędź do grafu widoczności.
				\item Zaktualizuj strukturę stanu dodając do niej odcinki incydentne leżące po stronie zgodnej z ruchem wskazówek zegara względem miotły i usuwając te leżące po stronie przeciwnej do ruchu wskazówek zegara.
			\end{itemize}
		\end{enumerate}
	
		\subsection{Struktura zdarzeń}
			Struktura zdarzeń została zaimplementowana przy użyciu listy. 
			
			\paragraph{Opis sortowania}
				Zbiór punktów najpierw dzielony jest na dwa podzbiory względem tego czy dany punkt leży poniżej czy powyżej aktualnie rozpatrywanego punktu $ p $ (w przypadku gdy leżą one na tej samej wysokości co punkt $ p $ podział następuje po wartości współrzędnej $ x $). Następnie każda część jest sortowana osobno wbudowaną funkcją $ sort $ (sortowanie stabilne) dwukrotnie. W pierwszym sortowaniu kluczem jest odległość od punktu zaczepienia, w drugim stosowany jest wyznacznik aby posortować punkty zgodnie z ruchem wskazówek zegara.
			
			\vspace{\baselineskip} % newline
			Ostatecznie otrzymujemy posortowaną strukturę zdarzeń gdzie punkty występują w kolejności zgodnie z ruchem wskazówek zegara, a punkty współliniowe ułożone są od najbliższych do najdalszych.
			
			\paragraph{Złożoność}
				Struktura zdarzeń nie jest modyfikowana w trakcie działania algorytmu, stąd jej wkład do złożoności jest ograniczony przez sortowanie i wynosi $ O(n \cdot log(n)) $.
		
		\subsection{Struktura stanu}
			Struktura stanu przechowuje krawędzie. Wykorzystuje posortowaną listę (opartą na drzewie binarnym) \lstinline|SortedList| z biblioteki \lstinline|sortedcontainers|. Gwarantuje ona czas operacji rzędu $ O(log(n)) $.
			
			\paragraph{Opis działania}
				Gdy miotła natrafia na punkt, brane są pod uwagę krawędzie incydentne z nim. Jeżeli krawędź została już zamieciona, jest ona usuwana ze struktury. Gdy nie została jeszcze zamieciona, jest dodawana do niej. Test po której stronie znajduje się krawędź jest przeprowadzany przy użyciu wyznacznika.
				
			\paragraph{Porządek elementów}
				Elementy są utrzymywane w porządku rosnącym względem ich odległości od punktu zaczepienia miotły. Jeżeli oba odcinki zaczynają lub kończą się w tym samym punkcie używamy wyznacznika.
				
			\paragraph{Złożoność}
				Na strukturze stanu wykonamy maksymalnie wykonamy $ n $ operacji wstawiania i usuwania, co daje nam czas działania rzędu $ O(n \cdot log(n)) $.
		
		\pagebreak
		\subsection{Sprawdzanie widoczności}
			
			Zgodnie z opisem algorytmu w punkcie \ref{sssec:algo} sprawdzanie widoczności będziemy wykonywać dla zadanych punktów $ p $ i $ w $. W celu określenia, czy punkt $ w $ jest widoczny będziemy analizować strukturę stanu miotły. W niektórych przypadkach potrzebujemy dodatkowe informacje o punkcie z struktury zdarzeń, który był rozważany poprzednio. Oznaczamy ten punkt jako $ prev\_w $.
			
			\vspace{\baselineskip} % newline
			Gdy zaczynamy rozpatrywać punkt, testujemy czy jest on widziany z punktu zaczepienia miotły.  W algorytmie spotykamy następujące przypadki:
			
			\paragraph{Przypadek 1}
			Punkty $ p $ i $ w $ należą do tej samej figury i wektor $ \overrightarrow{wp} $ jest skierowany do wnętrza wielokąta.
			
			
			\vspace{\baselineskip} % newline
			W tym przypadku punkt $ w $ nie jest widoczny, ponieważ linia łącząca $ p $ i $ w $ przechodziłaby przez wnętrze figury.
			
			\vspace{\baselineskip} % newline
			Przypadek musimy sprawdzić jako pierwszy, ponieważ taka sytuacja spełnia również przypadek 2 lub 3 i wtedy punkt $ w $ zostałby błędnie oznaczony jako widoczny.
			
			
			\img{przypadek1.jpg}{Przykład dla przypadku 1}
			
			Jak widać na przykładzie wektor $ \overrightarrow{wp} $ jest skierowany do wnętrza wielokąta, zatem pomimo że miotła nie jest przecinana przez żaden odcinek punkt $ w $ nie jest widoczny.
			
			\pagebreak
			\paragraph{Przypadek 2}
			Nie zaszedł przypadek 1 i struktura stanu miotły jest pusta.
			
			\vspace{\baselineskip} % newline
			Struktura stanu jest pusta zatem na linii z $ p $ do $ w $ nie może pojawić się żadna przeszkoda. Punkt $ w $ jest zatem widoczny.
			
			\img{przypadek2.jpg}{Przykład dla przypadku 2}
			
			\vspace{\baselineskip}

			\paragraph{Przypadek 3}
			Nie zaszedł żaden z powyższych przypadków i punkty $ p $, $ prev\_w $, $ w $ nie są współliniowe:
			
			\vspace{\baselineskip} % newline
			Jest to przypadek, który standardowo będzie występował najczęściej. Aby stwierdzić czy punkt $ w $ jest widoczny należy sprawdzić, czy punkt przecięcia odcinka znajdującego się na samym początku struktury stanu (czyli najbliżej $ p $) z miotłą leży bliżej, czy dalej od $ p $ niż $ w $. Inaczej mówiąc sprawdzamy, czy odcinek $ pw $ przecina odcinek z struktury stanu. Jeżeli tak jest to oznacza, że pomiędzy $ p $ i $ w $ znajduje się przeszkoda i punkt $ w $ jest niewidoczny. W przeciwnym wypadku punkt ten jest widoczny.
			
			\img{przypadek3.1.jpg}{Przykład 1 dla przypadku 3}
			
			Pierwszy odcinek z struktury stanu to $ DE $, przecina się on z odcinkiem $ pw $, zatem punkt $ w $ nie jest widoczny.
			
			\img{przypadek3.2.jpg}{Przykład 2 dla przypadku 3}
			
			Pierwszy odcinek z struktury stanu to $ KJ $, nie przecina się on z odcinkiem $ pw $, zatem punkt $ w $ jest widoczny.
			
			\paragraph{Przypadek 4}
			Nie zaszedł żaden z powyższych przypadków, punkty $ p $, $ prev\_w $, $ w $ są współliniowe oraz punkt $ prev\_w $ nie był widoczny:
			
			\vspace{\baselineskip} % newline
			Współliniowe punkty w strukturze zdarzeń, są ustawione w kolejności rosnącej odległości od $ p $, zatem jeżeli $ prev\_w $ był niewidoczny, to $ w $ także nie jest widoczny.
			
			\img{przypadek4.jpg}{Przykład dla przypadku 4}
			
			Punkt $ E $ był rozważany bezpośrednio przed punktem $ w $ i jest on niewidoczny. Dodatkowo punkty $ p $, $ E $, $ w $ są współliniowe. Wynika z tego, że punkt $ w $ również nie jest widoczny.
			
			\paragraph{Przypadek 5}
			Nie zaszedł żaden z powyższych przypadków, punkty $ p $, $ prev\_w $, $ w $ są współliniowe oraz punkt $ prev\_w $ był widoczny:
			
			\vspace{\baselineskip} % newline
			Przypadek ten jest podobny do przypadku 3, jednak musimy rozważyć go osobno, ponieważ po przetworzeniu punktu $ prev\_w $ na początku struktury stanu mogły zostać umieszczone odcinki, które nie przecinają miotły. Aby w tym przypadku stwierdzić, czy punkt $ w $ jest widoczny sprawdzamy, czy w strukturze stanu znajduje się odcinek, który przecina $ w\,prev\_w $. 
			
			\vspace{\baselineskip} % newline
			Gdy nie znajdziemy w strukturze stanu takiego odcinka, dodatkowo musimy rozważyć jeszcze jeden szczególny pod-przypadek, podobny do przypadku 1. Zachodzi on, gdy $ prev\_w $ i $ w $ należą do tej samej figury i wektor $ \overrightarrow{w\,prev\_w} $ jest skierowany do wnętrza wielokąta.
			
			\vspace{\baselineskip} % newline
			Jeśli zachodzi którakolwiek z opisanych powyżej sytuacji, to $ w $ jest niewidoczny. W przeciwnym wypadku punkt ten jest widoczny.
			
			
			\img{przypadek5.1.jpg}{Przykład 1 dla przypadku 5}
			
			Punkt $ D $ był rozważany bezpośrednio przed punktem $ w $ i jest widoczny. Punkty $ p $, $ D $, $ w $ są współliniowe. W strukturze stanu znajduje się odcinek $IJ$, który przecina $Dw$. Wynika z tego, że $ w $ jest niewidoczny.
			
			\img{przypadek5.2.jpg}{Przykład 2 dla przypadku 5}
			
			Punkt $ D $ był rozważany bezpośrednio przed $ w $ i jest widoczny, $ p $, $ D $, $ w $ są współliniowe. W strukturze stanu nie znajduje się odcinek, który przecina $Dw$, ale $ w $ i $ D $ należą do tej samej figury i wektor $ \overrightarrow{wD} $ jest skierowany do wnętrza wielokąta, zatem $ w $ jest niewidoczny.
			
			\img{przypadek5.3.jpg}{Przykład 3 dla przypadku 5}
			
			Punkt $ G $ był rozważany bezpośrednio przed punktem $ w $ i jest widoczny. Punkty $ p $, $ G $, $ w $ są współliniowe. W strukturze stanu nie znajduje się odcinek, który przecina $Gw$. Punkty $ w $ i $ G $ należą do różnych figur. Wynika z tego, że $ w $ jest widoczny.
			
			\img{przypadek5.4.jpg}{Przykład 4 dla przypadku 5}
			
			Punkt $ F $ był rozważany bezpośrednio przed punktem $ w $ i jest widoczny. Punkty $ p $, $ F $, $ w $ są współliniowe. W strukturze stanu nie znajduje się odcinek, który przecina $Fw$. Punkty $ w $ i $ F $ należą do tej samej figury, ale wektor $ \overrightarrow{wF} $ nie jest skierowany do wnętrza wielokąta, zatem $ w $ jest widoczny.
			
			\vspace{\baselineskip} % newline
			\vspace{\baselineskip} % newline
			\vspace{\baselineskip} % newline
			
			
			%\img{przyp1.jpg}{Przypadek graniczny}
			%	W tym przypadku najpierw dodany do struktury zostaje punkt $ A $. Następnie rozpatrujemy widoczność punktu $ E $. Algorytm najpierw sprawdza czy punkt poprzednio rozpatrywany należy do tej samej figury. Jeżeli tak, sprawdza czy odcinek $ AE $ przecina wnętrze figury. Jeżeli nie, punkt jest widoczny. Dopuszczamy widoczność wzdłuż krawędzi.
				
			%\img{przyp2.jpg}{Przypadek graniczny}
			%	Gdy punkt $ F $ leży ponad miotłą odcinek $ AE $ przecina wnętrze figury więc punkt $ E $ nie jest widoczny.
			
			%\img{przyp3.jpg}{Przypadek graniczny}
			%	Tutaj punkt $ D $ będzie widoczny. Sprawdzamy czy punkt $ A $ jest widoczny. Jeżeli tak, to wyszukujemy w strukturze stanu czy pomiędzy punktami $ A $ i $ D $ znajduje się jakakolwiek krawędź. Jeżeli nie, rozpatrywany punkt jest widoczny.
				
			%\img{przyp4.jpg}{Przypadek graniczny}
			%	W tym przykładzie punkt $ G $ nie będzie widoczny ponieważ pomiędzy punktem $ A $ (który jest widoczny) a punktem $ G $ w strukturze stanu znajdują się dwie krawędzie.
				
			%\img{przyp5.jpg}{Przypadek graniczny}
			%	W tym przykładzie punkt $ D $ nie ma współliniowych poprzedników w strukturze zdarzeń. W takim przypadku wyszukujemy pierwszy odcinek w strukturze stanu i jeżeli taki istnieje i przecina miotłę to punkt nie będzie widoczny. Szukanie przeszkód przecinających odcinek $ DG $ w tym przypadku nie ma sensu. $ G $ również jest niewidoczny.
			
			%\img{przyp6.jpg}{Przypadek graniczny}
			%	Tutaj odcinek $ AB $ pokrywa się z miotłą. Dopuszczamy widoczność wzdłuż krawędzi. Punkt $ D $ będzie widoczny.
		
		\subsection{Kroki działania algorytmu} \label{sssec:steps}
			Punkt $ A $ jest punktem zaczepienia miotły. Miotła jest reprezentowana jako czerwony odcinek z $ A $ do obecnie przetwarzanego punktu.
			
			\imgcustomsize{visibility-step000-001.jpg}{Przykład działania}{0.35}
			Algorytm napotyka pierwsze zdarzenie.  Punkt $ D $ jest widoczny, ponieważ struktura stanu jest pusta. Odcinki incydentne $ DB $ i $ DC $ dodawane są do struktury stanu (w tej kolejności).
			
			\imgcustomsize{visibility-step000-002.jpg}{Przykład działania}{0.35}
			$ DB $ jest usuwane ze struktury, $ BC $ zostaje dodane. Punkt $ B $ jest widoczny, ponieważ nie istnieje krawędź przecinająca miotłę pomiędzy $ A $ i $ B $.
			
			\imgcustomsize{visibility-step000-003.jpg}{Przykład działania}{0.35}
			Dodajemy do struktury $ GF $ i $ GE $. Punkt $ G $ nie jest widoczny, bo zasłania go krawędź $ BC $ (pierwsza przecinająca miotłę).
			
			\imgcustomsize{visibility-step000-004.jpg}{Przykład działania}{0.35}
			Usuwamy $ GE $, dodajemy $ EF $. Analogicznie do $ G $, punkt $ E $ nie jest widoczny.
			
			\imgcustomsize{visibility-step000-005.jpg}{Przykład działania}{0.35}
			Usuwamy $ EF $ i $ GF $. Punkt $ F $ niewidoczny.
			
			\imgcustomsize{visibility-step000-006.jpg}{Przykład działania}{0.35}
			Usuwamy $ BC $ i $ DC $. Punkt $ C $ widoczny.
			
			\imgcustomsize{visibility-sum-001.jpg}{Wynik działania}{0.35}
			Tak wygląda wynik wykonania powyższych kroków. Następnie wybierany jest nowy punkt zaczepienia i algorytm powtarza kroki wykonane powyżej.
			
			\vspace{\baselineskip} % newline
			Kolorem czerwonym oznaczone są figury (przeszkody), krawędzie dodane do grafu widoczności zaznaczone zostały przerywanymi liniami. 
			
			\imgsidebyside{visibility-sum-002.jpg}{Miotła zaczepiona w $ B $}{visibility-sum-003.jpg}{Miotła zaczepiona w $ C $}
			\imgsidebyside{visibility-sum-004.jpg}{Miotła zaczepiona w $ D $}{visibility-sum-005.jpg}{Miotła zaczepiona w $ E $}
			\imgsidebyside{visibility-sum-006.jpg}{Miotła zaczepiona w $ F $}{visibility-sum-007.jpg}{Miotła zaczepiona w $ G $}
			
			Tak wygląda ostatecznie graf widoczności dla tego przypadku.
		
		\subsection{Reprezentacja grafu}

		Graf jest zwracany przez funkcję $ create\_visibility\_graph $ jako obiekt klasy $ VisibilityGraph $. Wewnątrz klasy graf jest reprezentowany jako słownik słowników. Dla każdego wierzchołka (reprezentowanego jako obiekt klasy $ Point $) przechowywany jest słownik jego sąsiadów (widocznych punktów), gdzie kluczem jest sąsiadujący wierzchołek (również reprezentowany jako obiekt klasy $ Point $), a wartością jest odległość tych dwóch punktów na płaszczyźnie. Graf w takiej postaci zwracany jest przez metodę $ get\_graph $.
		
		\vspace{\baselineskip} % newline
		Klasa posiada również metodę $ get\_lines $, która zwraca listę krawędzi grafu w postaci par punktów [$ pierwszy $ $ punkt $, $ drugi $ $ punkt $]. Punkty te są w tej liście reprezentowane jako dwuelementowe krotki postaci (współrzędna $ x $ punktu, współrzędna $ y $ punktu).
		
		\vspace{\baselineskip} % newline
		Poniżej w Tablicy 1 przedstawiony został graf widoczności dla przypadku opisywanego w punkcie \ref{sssec:steps}, w postaci listy krawędzi.
		
		\tab{res/results.txt}{Reprezentacja grafu widoczności poprzez pary punktów tworzące listę krawędzi}
		
		\subsection{Testy}
			\imgcustomsize{test1.jpg}{Przypadek testowy}{0.75}
				Wierzchołek $ F $ jest widoczny, ponieważ $ B $ jest widoczny i miotły nie przecinają żadne krawędzie na odcinku $ BF $ oraz wektor $ \overrightarrow{FB} $ nie jest skierowany do wnętrza figury.
			
			\imgcustomsize{test2.jpg}{Przypadek testowy}{0.75}
				Tutaj wektor $ FB $ jest skierowany do wnętrza figury więc punkt $ F $ jest niewidoczny (przerywana linia ma za zadanie tylko podkreślić fakt, że punkty $ B $ i $ F $ są współliniowe).
			
			\imgcustomsize{test3.jpg}{Przypadek testowy}{0.75}
				Punkt $ E $ jest widoczny, ponieważ $ B $ jest widoczny i należą do różnych figur oraz pomiędzy nimi nie ma żadnej krawędzi przecinającej odcinek $ BE $.
			
			\imgcustomsize{test4.jpg}{Przypadek testowy}{0.75}
				Tutaj analogicznie jak wyżej, ale $ BH $ jest przecięty dwoma odcinkami wiec $ H $ nie będzie widoczny.
			
			\imgcustomsize{test5.jpg}{Przypadek testowy}{0.75}
				Jeżeli $ E $ jest niewidoczny to nie ma sensu badać widoczności $ H $, ponieważ jest on dalej od punktu zaczepienia miotły (punktu $ A $).
			
			\imgcustomsize{test6.jpg}{Przypadek testowy}{0.75}
				Wierzchołek $ E $ widoczny, gdy miotła pokrywa się z krawędzią $ BC $.
				
			\imgcustomsize{test7.jpg}{Przypadek testowy}{1}
			\imgcustomsize{test8.jpg}{Przypadek testowy}{1}
			\imgcustomsize{test9.jpg}{Przypadek testowy}{1}
			\imgcustomsize{test10.jpg}{Przypadek testowy}{1}
		
		\subsection{Testy losowe}
			% TODO opis testy losowe
			\imgcustomsize{test-losowe1.jpg}{Przypadek testowy}{1}
			\imgcustomsize{test-losowe2.jpg}{Przypadek testowy}{1}
			\imgcustomsize{test-losowe3.jpg}{Przypadek testowy}{1}
			\imgcustomsize{test-losowe4.jpg}{Przypadek testowy}{1}
		
		\subsection{Testy znajdowania najkrótszej ścieżki}
			\imgcustomsize{p1.jpg}{Przypadek testowy}{0.75}
			\imgcustomsize{p2.jpg}{Przypadek testowy}{0.75}
		
	\section{Bibliografia}
		Algorytm został przygotowany na podstawie książki Marka de Berga pt. \textit{Geometria Obliczeniowa - Algorytmy i Zastosowania}.

\end{document}
